% =====================================================================
% LaTeX INSERT PACK — ICCEREC 2026 camera-ready revision
% Paper: Adaptive Trust-Aware Access Control for Zero Trust Security
%        in OpenDaylight-Based Software-Defined Networks
%
% This file contains verified numbers produced by:
%   vm-controller/PDP/sensitivity_analysis.py
% Run there (VM1):  cd ~/PDP && python3 sensitivity_analysis.py
% Raw evidence:     sensitivity_results.md / sensitivity_results.csv
%
% Copy the blocks below into your manuscript where indicated.
% =====================================================================


% ---------------------------------------------------------------------
% (R1e) Proper mathematical typesetting of the trust-score formula
% Replace the plain "T = 0.5R+0.3C+0.2B" text with numbered equations.
% ---------------------------------------------------------------------
The trust score $T$ of an access request is a weighted combination of three
normalised factors (each on a $0$--$100$ scale):

\begin{equation}
  T = w_R R + w_C C + w_B B,
  \label{eq:trust}
\end{equation}

\noindent where $R$, $C$, and $B$ denote role/identity, contextual, and
behavioural trust, respectively, and $w_R=0.5$, $w_C=0.3$, $w_B=0.2$
($w_R+w_C+w_B=1$). The contextual and behavioural scores are obtained by
point deductions:

\begin{align}
  C &= \max\!\big(0,\; 100 - p_{ip}\,\mathbb{1}_{ip\notin S}
        - p_{h}\,\mathbb{1}_{h\notin H}
        - p_{m}\,\mathbb{1}_{\text{mac}\neq\text{mac}^\ast}\big),
  \label{eq:context}\\
  B &= \max\!\big(0,\; 100 - \min(n\,p_{f},\, p_{f}^{\max})\big),
  \label{eq:behaviour}
\end{align}

\noindent with $p_{ip}=40$, $p_h=30$, $p_m=50$, $p_f=15$, $p_f^{\max}=60$,
and $n$ the number of consecutive failed logins. The resulting score is
mapped to three access levels:

\begin{equation}
  \text{tier}(T) =
  \begin{cases}
    \text{FULL},    & T \ge 70,\\
    \text{LIMITED}, & 40 \le T < 70,\\
    \text{DENIED},  & T < 40.
  \end{cases}
  \label{eq:tier}
\end{equation}


% ---------------------------------------------------------------------
% (R1b) Figure 1 caption — define all H1-H4 and S1-S4 explicitly
% ---------------------------------------------------------------------
\begin{figure}[t]
  \centering
  \includegraphics[width=\linewidth]{fig/architecture.pdf}
  \caption{System architecture of the proposed ODL-ZTA framework. The
  control/policy plane runs on VM1 (OpenDaylight controller and the Flask
  Policy Decision Point), and the data plane runs on VM2 (Mininet with four
  Open vSwitch instances). Hosts $h_1$ (research, $10.0.0.1$), $h_2$ (server,
  $10.0.0.2$), $h_3$ (IoT, $10.0.0.3$), and $h_4$ (guest, $10.0.0.4$) each
  attach to switch port~1 of switches $s_1$--$s_4$; the switches form a ring
  $s_1$--$s_2$--$s_3$--$s_4$--$s_1$.}
  \label{fig:architecture}
\end{figure}


% ---------------------------------------------------------------------
% (R1a/R3-1) Scenario validation table with R/C/B provenance
% Matches Table II of the current manuscript.
% ---------------------------------------------------------------------
\begin{table}[t]
  \caption{Evaluated scenarios: provenance of the trust factors and the
  resulting access level.}
  \label{tab:scenarios}
  \centering
  \begin{tabular}{llcccr}
    \hline
    \textbf{Expected tier} & \textbf{Account (role)} & $R$ & $C$ & $B$ & $T$ \\
    \hline
    FULL    & ratih (research) & 80 & 100 & 100 & 90.0 \\
    LIMITED & ratih (research) & 80 &  50 &  70 & 69.0 \\
    DENIED  & bima (guest)     & 30 &  50 &  40 & 38.0 \\
    \hline
  \end{tabular}
  \\[2pt]
  \footnotesize $C=50$ follows from an IP/MAC binding mismatch ($-50$);
  $B=70$ and $B=40$ follow from 2 and 4 failed logins, respectively.
\end{table}


% ---------------------------------------------------------------------
% (R1a/R2/R3-1) Weight sensitivity table — the core reviewer answer
% ---------------------------------------------------------------------
\begin{table}[t]
  \caption{Sensitivity of access-level outcomes to the trust-score weights
  (thresholds fixed at 70/40). The alternative raised by the reviewers
  (Alt-1: $0.45/0.25/0.30$) preserves all three outcomes.}
  \label{tab:weight-sensitivity}
  \centering
  \begin{tabular}{lcccccc}
    \hline
    \textbf{Configuration} & $w_R$ & $w_C$ & $w_B$
      & \textbf{FULL} & \textbf{LIMITED} & \textbf{DENIED} \\
    \hline
    Baseline           & 0.50 & 0.30 & 0.20 & 90.0 (FULL)    & 69.0 (LIMITED) & 38.0 (DENIED) \\
    Alt-1 (flatter)    & 0.45 & 0.25 & 0.30 & 91.0 (FULL)    & 69.5 (LIMITED) & 38.0 (DENIED) \\
    Alt-2 (context)    & 0.40 & 0.40 & 0.20 & 92.0 (FULL)    & 66.0 (LIMITED) & 40.0 (LIMITED) \\
    Alt-3 (behaviour)  & 0.40 & 0.20 & 0.40 & 92.0 (FULL)    & 70.0 (FULL)    & 38.0 (DENIED) \\
    Alt-4 (identity)   & 0.60 & 0.30 & 0.10 & 88.0 (FULL)    & 70.0 (FULL)    & 37.0 (DENIED) \\
    Alt-5 (uniform)    & 0.33 & 0.33 & 0.33 & 93.3 (FULL)    & 66.7 (LIMITED) & 40.0 (LIMITED) \\
    \hline
  \end{tabular}
\end{table}


% ---------------------------------------------------------------------
% (R1a) Threshold sensitivity table
% ---------------------------------------------------------------------
\begin{table}[t]
  \caption{Sensitivity of access-level outcomes to the tier thresholds
  (weights fixed at $0.5/0.3/0.2$).}
  \label{tab:threshold-sensitivity}
  \centering
  \begin{tabular}{lccccc}
    \hline
    \textbf{Thresholds} & \textbf{FULL case} & \textbf{LIMITED case} & \textbf{DENIED case} \\
    \hline
    70 / 40 (baseline) & 90.0 $\rightarrow$ FULL    & 69.0 $\rightarrow$ LIMITED & 38.0 $\rightarrow$ DENIED \\
    65 / 35 (looser)   & 90.0 $\rightarrow$ FULL    & 69.0 $\rightarrow$ FULL    & 38.0 $\rightarrow$ LIMITED \\
    75 / 45 (stricter) & 90.0 $\rightarrow$ FULL    & 69.0 $\rightarrow$ LIMITED & 38.0 $\rightarrow$ DENIED \\
    70 / 50            & 90.0 $\rightarrow$ FULL    & 69.0 $\rightarrow$ LIMITED & 38.0 $\rightarrow$ DENIED \\
    60 / 30 (lower)    & 90.0 $\rightarrow$ FULL    & 69.0 $\rightarrow$ FULL    & 38.0 $\rightarrow$ LIMITED \\
    \hline
  \end{tabular}
\end{table}


% ---------------------------------------------------------------------
% (R1a) Threshold-placement / robustness table + paragraph
% ---------------------------------------------------------------------
\begin{table}[t]
  \caption{Robustness of the FULL/LIMITED/DENIED mapping of the three
  evaluated scenarios over sampled weight assignments.}
  \label{tab:robustness}
  \centering
  \begin{tabular}{lc}
    \hline
    \textbf{Weight domain} & \textbf{Mapping preserved} \\
    \hline
    Full simplex (step 0.05)          & 27/231 (11.7\%) \\
    Design-consistent $w_R\ge w_C\ge w_B$ & 79/154 (51.3\%) \\
    Near baseline ($\pm 0.05$)        & 21/25 (84.0\%) \\
    \hline
  \end{tabular}
\end{table}

The thresholds are not arbitrary. The set of scores reachable from valid
$R/C/B$ combinations contains neither exactly $70$ nor exactly $40$: the
nearest values are $69.0$ and $70.5$ around the FULL boundary, and $39.0$ and
$41.0$ around the LIMITED boundary. Consequently, every threshold within
$(69.0, 70.5]$ and $(39.0, 41.0]$ yields an identical tiering, so the chosen
values of $70$ and $40$ are invariant to small perturbations.


% ---------------------------------------------------------------------
% (R1a/R3-1) Suggested paragraph for the Results/Discussion section
% ---------------------------------------------------------------------
A sensitivity analysis was performed to address the choice of weights and
thresholds. Table~\ref{tab:weight-sensitivity} shows that the alternative
weighting $w_R,w_C,w_B = 0.45,0.25,0.30$ raised by the reviewers produces the
same access-level decisions as the proposed $0.5,0.3,0.2$ (scores $91.0$,
$69.5$, and $38.0$ for the FULL, LIMITED, and DENIED scenarios, respectively).
The mapping changes only when identity is under-weighted (Alt-2 and Alt-5,
where the guest scenario reaches exactly $40$ and escapes denial) or when a
single non-identity factor dominates (Alt-3 and Alt-4, where a session with a
compromised context reaches exactly $70$ and is granted FULL access). This
confirms that the proposed weighting is appropriate specifically because it
keeps identity dominant without letting any single factor override the
contextual signal.

\begin{itemize}
  \item Table~\ref{tab:threshold-sensitivity}: the tiering is unchanged under
        moderate threshold shifts ($75/45$ and $70/50$); only substantial
        loosening ($65/35$ or $60/30$) alters the outcomes.
  \item Table~\ref{tab:robustness}: near the chosen values ($\pm 0.05$) the
        mapping is preserved for $84.0\%$ of sampled weight assignments, and
        for the identity-dominant region in $51.3\%$ of cases. The mapping
        degrades only for weightings far from that region; the affected cases
        are precisely those placed on a tier boundary, which reflects the
        intended boundary behaviour of the adaptive policy.
\end{itemize}

These results are obtained by executing
\path{sensitivity_analysis.py}, which evaluates Eq.~\eqref{eq:trust} directly
over the reachable $R/C/B$ space and is fully reproducible without the data
plane.


% =====================================================================
% (R3-4) Performance metrics — measured on the live testbed
% Source: perf_collect.py  (20 login/logout cycles + 4 concurrent sessions)
% =====================================================================
\begin{table}[t]
  \caption{Policy Decision Point latency during 20 full access sessions
  (FULL tier, identity-dominant policy).}
  \label{tab:latency}
  \centering
  \begin{tabular}{lcccc}
    \hline
    \textbf{Operation} & \textbf{Mean} & \textbf{p50} & \textbf{p95} & \textbf{Max} \\
    \hline
    Provisioning (login)  & 143.5 ms & 139.6 ms & 197.1 ms & 226.5 ms \\
    Revocation (logout)   &  94.3 ms &  90.8 ms & 138.4 ms & 167.1 ms \\
    \hline
  \end{tabular}
\end{table}

\begin{table}[t]
  \caption{OpenDaylight controller load. Each full session installs 16
  OpenFlow rules; concurrent revocation removes the rules of four sessions.}
  \label{tab:controller-load}
  \centering
  \begin{tabular}{ll}
    \hline
    \textbf{Metric} & \textbf{Value} \\
    \hline
    RESTCONF PUT per rule (average)        & 7.84 ms \\
    RESTCONF PUT per rule (p95)            & 13.06 ms \\
    Concurrent sessions revoked            & 4 \\
    Flows revoked concurrently             & 64 \\
    Total concurrent revocation time       & 248.1 ms \\
    Controller CPU during session churn    & 140.8\% \\
    \hline
  \end{tabular}
\end{table}

The performance impact of policy provisioning and revocation was measured on
the live testbed. Across 20 full access sessions, the PDP provisioned the
per-session OpenFlow rules in 143.5~ms on average ($p95 = 197.1$~ms), and
revoked them in 94.3~ms on average ($p95 = 138.4$~ms). Each full session
installs 16 rules, so the per-rule RESTCONF latency averages 7.84~ms. During a
concurrent revocation of four sessions (64 rules), the total revocation time
was 248.1~ms and the OpenDaylight controller consumed 140.8\% CPU over the
$\approx 4.9$~s session-churn window, confirming that enforcement is viable
without specialised hardware while indicating that controller load grows with
the number of concurrently enforced sessions.


% ---------------------------------------------------------------------
% (R3-4) Data-plane throughput and per-tier reachability
% ---------------------------------------------------------------------
\begin{table}[t]
  \caption{Data-plane reachability and throughput per access state.}
  \label{tab:throughput}
  \centering
  \begin{tabular}{lccc}
    \hline
    \textbf{State} & \textbf{server:9000} & \textbf{server:8080} & \textbf{IoT:80} \\
    \hline
    Default-deny (pre-login) & blocked & blocked & blocked \\
    FULL                     & 71.0 Gbit/s & HTTP 200 & HTTP 200 \\
    LIMITED                  & blocked & blocked & HTTP 200 \\
    Post-logout              & blocked & blocked & blocked \\
    \hline
  \end{tabular}
\end{table}

The data-plane effect of each access tier was verified with \texttt{iperf3} on
the server port 9000 and with HTTP requests to the authorised resources. Under
FULL access, the session reached \textbf{71.0~Gbit/s} (sender; 70.6~Gbit/s
receiver) over the Open vSwitch path, compared with approximately 85~Gbit/s for
the loopback baseline of the same host. Under LIMITED access, both server ports
were blocked while the IoT service on port 80 still returned HTTP~200. Under
default-deny (before authentication) and after logout, all three destinations
were blocked. These results show that the tier decisions of the PDP are enforced
as resource-specific reachability at the data plane, not only as a functional
state.


% ---------------------------------------------------------------------
% (R3-2) Why a deterministic formulation instead of ML/AI
% ---------------------------------------------------------------------
\subsection{Rationale for a Deterministic Trust Formulation}
\label{sec:why-not-ml}
The proposed mechanism uses explicit point deductions rather than a learned
model. This is a deliberate design choice for three reasons. First,
network-policy enforcement requires decisions that are deterministic and
auditable: every tier assignment in this work can be traced to a specific factor
and penalty, whereas an LSTM- or CNN-based classifier used for anomaly detection
\cite{guo2023intelligent} produces scores whose provenance is difficult to
attach to a concrete flow rule. Second, the mechanism operates without labelled
training data, which is scarce for this access-control context, and it runs
alongside the controller on commodity hardware without a training pipeline or
GPU. Third, learned traffic-anomaly detection and transparent policy decisions
are complementary rather than competing: an ML model could supply a richer
behavioural signal that replaces or augments the $B$ factor while the
transparent aggregation in Eq.~\eqref{eq:trust} is retained. The main limitation
of the point-deduction approach is that it cannot capture complex non-linear
adversarial patterns; enriching the behavioural factor with learned anomaly
scores is therefore identified as future work.


% ---------------------------------------------------------------------
% (R3-3) Scope of topological validation (justification, option B)
% ---------------------------------------------------------------------
\subsection{Scope of Topological Validation}
\label{sec:topology-scope}
The mechanism is topology-agnostic: the PDP derives each path with a
breadth-first search over an adjacency graph and installs one flow per switch
along that path, so no component assumes a ring. The four-switch ring is used
as a controlled minimal topology that contains a cycle, which exercises both the
default-deny baseline and loop handling while keeping the trust-to-policy
mapping isolated from unrelated effects. The number of switches does not affect
the trust computation, which depends only on the $R$, $C$, and $B$ factors, nor
the RESTCONF provisioning procedure. The per-session deployment cost is bounded
by the path length and the number of authorised ports, scaling linearly with the
network size. Consequently, the evaluation establishes functional correctness
and the correctness of the trust-to-enforcement mapping, not enterprise-scale
viability; validation on larger reference topologies such as Abilene or GEANT is
left as future work. One implementation requirement that emerged during testing
is that the port mapping used by the PDP must exactly match the topology built
in the data plane, since a mismatch silently misdirects return traffic without
raising an error.


% =====================================================================
% (Opsional) Paragraf batasan untuk Conclusion
% =====================================================================
% The evaluation does not establish enterprise-scale viability or resistance to
% attack; it validates the functional mapping from trust conditions to
% resource-specific enforcement and reports measured controller-side and
% data-plane performance for a four-switch ring topology.

